We present \system, a sign-transcription architecture designed to connect a
lightweight visual model to a pretrained language model through its native token
interface.  The original \system{} proposal imitated continuous language-model
embeddings.  We instead predict explicit \Gemma{} token IDs: this removes the
ambiguity of decoding an approximate vector, permits exact sequence evaluation,
and gives a downstream \Gemma{} model a standard textual input that it can
punctuate, normalize, or---when context permits---correct.  A spatial--temporal
keypoint encoder and autoregressive decoder are trained on synthetic streams of
two to eight LSA64 signs and evaluated with signer-held-out splits.  On three
development signers, the selected model improves strict token-sequence accuracy
from 1.8--2.3\% for the initial autoregressive decoder to 16.9--19.8\% while
preserving the exact \Gemma{} IDs.  Prospectively registered evaluation on two
additional signers obtains \FoldSevenExact\% and \FoldEightExact\% strict
accuracy.  We also audit the complete token-to-language interface.  \Gemma{} 3n
preserves 95\% of clean isolated-label inputs, but recovers no exact sequence at
any tested nonzero substitution rate and slightly decreases chrF.  This negative
result identifies a current boundary: random isolated signs contain little
sentence context from which to infer a corrupted word. These findings are limited
to this seed-23 evaluation of synthetic, two-to-eight-sign LSA64 streams and two
prospectively held-out signers. The isolated-label audit does not establish
correction in natural sentences.
